\documentclass[10pt,a4paper]{amsart}
\usepackage[margin=19mm]{geometry}
\usepackage{amsmath,amssymb,mathtools}
\usepackage[scr=rsfs,cal=euler]{mathalfa}
\usepackage{xfrac}
\usepackage[unicode,hidelinks,bookmarks=false]{hyperref}
\newcommand{\ie}{i.e.\ }
\newcommand{\cf}{cf.\ }
\newcommand{\resp}{respectively\ }
\newcommand{\Subsection}{\subsection}
\providecommand{\nobreakdash}{\nobreak\hbox{-}}
\setlength{\emergencystretch}{2em}
\multlinegap0pt
\usepackage{bm}
\usepackage{bbm}
\usepackage{dsfont}
\usepackage{xspace}
\usepackage{cancel}
\usepackage{mathtools}
\usepackage{amsthm}
\makeatletter
\@ifundefined{theorem}{\newtheorem{theorem}{Theorem}}{}
\makeatother
\title[Arithmetic invariants of torus links]{Arithmetic invariants of torus links}
\author{Anwesh Ray, Tanushree Shah}
\date{Source: arXiv:2511.03446v1 (CC BY 4.0)}
\begin{document}
\maketitle

\noindent\emph{Source and licence.} Arithmetic invariants of torus links. Version arXiv:2511.03446v1 (CC BY 4.0), distributed under the Creative Commons Attribution 4.0 International licence, as recorded in the arXiv licence metadata for this exact version and shown on the abstract page https://arxiv.org/abs/2511.03446v1. Full rights notice: licenses/arxiv-2511.03446-CC-BY-4.0.txt.\par\smallskip

\noindent\textit{Editorial scope.} Counted unit: the Theorem whose source label is \texttt{section 3 main thm 2} in the article (source0.tex lines 836--839, source label \texttt{section 3 main thm 2}), delivered together with the article's own complete original proof of it (source0.tex lines 840--928, one \texttt{proof} environment of 89 lines). No mathematics is written, repaired or re-derived: statement and proof are the article's own text line for line. Required context retained uncounted: none. Every definition, display and label of those lines is retained unchanged. Omitted from this package and not redistributed: 1--835, 929--1757. Nothing inside a retained excerpt is omitted or altered: every retained body is delivered exactly as the article's lines are written, so no omission receipt of any kind accompanies this package. Citations: the retained excerpts carry no citation command at all, so no bibliography entry of the article is transcribed and no reference is redistributed. Reference adaptation: none was needed and none was made; every reference inside the delivered unit resolves inside it, so no unresolved reference or citation is produced. Printed slips kept literal: no printed word, symbol, hypothesis or number of the counted statement or of its proof was altered. Layout and class: the article's own class is \texttt{amsart} with packages including geometry, amscd, amsmath, amsxtra, amsthm, amssymb, stmaryrd, xr, mathrsfs, mathtools, enumerate, commath, comment, tikz; the wrapper is the lane's pinned \texttt{amsart} editorial wrapper at 10pt on A4, which restates the theorem-like environments the retained text needs and adds the article's own macro definitions that the retained text needs (none), with extra wrapper packages (bm, bbm, dsfont, xspace, cancel, mathtools). The delivered numbering is the unit's own. Assets: no retained span contains a figure environment, a graphics-inclusion command, a PDF-page inclusion, a table or a TikZ picture; no image file is copied into this package, the package declares no asset dependency and no image file is redistributed with it. Licence: version arXiv:2511.03446v1 of this article is distributed under the Creative Commons Attribution 4.0 International licence, as recorded in the arXiv licence metadata for this exact version and shown on the abstract page https://arxiv.org/abs/2511.03446v1. A full rights notice is delivered at licenses/arxiv-2511.03446-CC-BY-4.0.txt. Characters: every retained excerpt is pure ASCII, so the delivered engine, XeLaTeX with its default Latin Modern text font, reports no missing character.
\begin{theorem}\label{section 3 main thm 2}
    With respect to notation above, 
\begin{equation}\label{Omega asymptotic 1}\#\Omega_{[a,b]}(X)=\frac{(b-a)}{4} X^4+O(X^3)\end{equation}\noindent and\begin{equation}\label{Omega asymptotic 2}\lim_{X\rightarrow \infty} \frac{\#\Omega_{[a,b]}(X)}{\#\Omega(X)}.\end{equation}
\end{theorem}

\begin{proof}
\par We write each cyclotomic factor as \(t^n-1=\prod_{r\mid n}\Phi_r(t)\). Comparing exponents of the cyclotomic polynomials \(\Phi_r\) in the expression for \(\Delta_{p,q}\) shows that for each \(r\ge1\)
\[
M_r= d \mathbf{1}_{r\mid L}-\mathbf{1}_{r\mid p}-\mathbf{1}_{r\mid q} + \mathbf{1}_{r=1},
\]
where \(\mathbf{1}_{P}\) denotes the indicator of the predicate \(P\). In particular:
\begin{itemize}
    \item If $r$ does not divide $L$ then $M_r=0$.
 \item If $r$ divides $L$ but does not divide $p$ or $q$, then $M_r=d$.
  \item If \(r>1\) divides \(d\) then
  \(M_r=(d-2)\).
  \item If \(r>1\) divides $p$ (resp. $q$) but does not divide $q$ (resp. $p$), then \(M_r=(d-1)\).
  \item For the linear factor \(r=1\) (i.e. \(t-1\)) the exponent is \(M_r=(d-1)\).
\end{itemize}
In particular every root of \(\Delta_{p,q}\) is a root of unity whose order divides \(L\), and the multiplicity of any given primitive \(r\)-th root of unity is bounded above by \(d\). 
\par For a pair \((p,q)\) let
\[R_{p,q}:=Ld+1-p-q=(p-1)(q-1)\] be the total number of roots of \(\Delta_{p,q}\) counted with multiplicity. Thus one has that
\[
\#\Omega(X):=\sum_{1\le p,q\le\lfloor X\rfloor} R_{p,q}=\sum_{1\le p,q\le\lfloor X\rfloor} (p-1)(q-1)=\frac{1}{4}\left(\lfloor X\rfloor(\lfloor X\rfloor-1)\right)^2=\frac{1}{4}X^4+O(X^3).
\]
Since every root is a root of unity of order dividing \(L\),
we can enumerate them as \(e^{2\pi i k/L}\) for \(k=1,\dots,L-1\) with certain multiplicities. Let $S_r$ denote the multiset of primitive $r$-th roots of unity, each counted with multiplicity $M_r$.  Then the full multiset of zeros of $\Delta_{p,q}$ (with multiplicity) is the disjoint union $\bigsqcup_{r\mid L} S_r$, and 
\[\sum_{r\mid L} M_r\varphi(r)
=R_{p,q}=(p-1)(q-1),
\]
where $\varphi$ is Euler's totient function.
\par For any interval $[a,b]\subset[0,1]$, define
\[
N_r([a,b]) := \#\{1\le k\le r : \gcd(k,r)=1,\ \tfrac{k}{r}\in[a,b]\}.
\]
Letting $\mathbf{1}_{[a,b]}$ be the indicator functor for containment in $[a,b]$ defined as follows:
\[\mathbf{1}_{[a,b]}(x):=\begin{cases}
    & 1\text{ if }x\in [a,b];\\
    & 0\text{ if }x\notin [a,b].
\end{cases}\]

One finds that 
\[\begin{split}N_r([a,b])=&\sum_{k=1}^r \mathbf{1}_{[a,b]}(k/r)\sum_{d|(k,r)} \mu(d)\\
=& \sum_{d|r} \mu(d)\sum_{k=1}^{r/d} \mathbf{1}_{[a,b]}(kd/r)\\
=& \sum_{d|r} \mu(d)\sum_{k=1}^{r/d} \mathbf{1}_{[a,b]}(k/(r/d))\\
=& \sum_{d|r} \mu(d)\sum_{k=1}^{r/d} \mathbf{1}_{[ra/d,rb/d]}(k)\\
=& (b-a)\sum_{d|r} \mu(d)\frac{r}{d}+O\left(\sum_{d|r} 1\right)\\
=& (b-a)\varphi(r) + O(d(r))
\end{split}\]
where $d(r)$ denotes the number of positive divisors of $r$. In particular, for every $\varepsilon>0$, there exists a constant $C_\varepsilon>0$
such that
\[
N_r([a,b]) = (b-a)\varphi(r) + \delta_r, \qquad |\delta_r|\le C_\varepsilon r^\varepsilon.
\]
\noindent Multiplying by the multiplicity $M_r$ of the factor $\Phi_r(t)$ in
$\Delta_{p,q}(t)$ yields
\[
\#\{\xi\in S_r\mid \theta(\xi)\in[a,b]\}
= M_r N_r([a,b])
= M_r\big((b-a)\varphi(r)+\delta_r\big),
\qquad |\delta_r|\ll d(r).
\]
\noindent Summing over all divisors $r\mid L$, we obtain
\begin{equation}\label{eq:local-distribution}
\begin{split}
& \sum_{r|L}\# \{\xi\in S_r\mid \theta(\xi)\in [a,b]\}\\
=&\sum_{r|L} M_r N_{p,q}([a,b])\\
=&(b-a)\sum_{r\mid L} M_r\varphi(r)
+\sum_{r\mid L} M_r\delta_r
=(b-a)R_{p,q}+E_{p,q},
\end{split}
\end{equation}
where the total error term is
\[
E_{p,q}=\sum_{r\mid L}M_r\delta_r\leq d\sum_{r|L} \delta_r\ll_{\epsilon} d L^{\epsilon}.
\]
\noindent Recall that $p=dp'$, $q=dq'$ with $(p',q')=1$. Thus writing $L=d p'q'$ with $p',q'\le X/d$, one finds that
\[\begin{split}
&\sum_{1\le p,q\le X} |E_{p,q}| \\
\ll_{\epsilon} &\sum_{d\le X} d^{1+\epsilon} \sum_{p',q'\le X/d} (p'q')^\epsilon\\
\leq &\sum_{d\le X} d^{1+\epsilon} \left(\frac{X}{d}\right)^{2+2\epsilon} \\
\leq & X^{2+2\epsilon}\sum_{d\le X} \frac{1}{d^{1+\epsilon}}
\ll_{\epsilon} X^{2+\epsilon}.
\end{split}
\]
Thus we find that 
\[\begin{split}
    \# \Omega_{[a,b]}(X)= &
    \sum_{1\leq p,q\leq X}\sum_{r|L(p,q)}\# \{\xi\in S_r\mid \theta(\xi)\in [a,b]\}\\
    = &(b-a) \sum_{1\leq p,q\leq X} R_{p,q}+\sum_{1\leq p,q\leq X}E_{p,q}\\
    = &(b-a)\#\Omega(X)+O_{\epsilon}(X^{2+\epsilon}).
\end{split}\]
Since $\#\Omega(X)=\frac{1}{4} X^4+O(X^3)$, \eqref{Omega asymptotic 1} and \eqref{Omega asymptotic 2} follow from the above.
\end{proof}

\end{document}
